\documentclass[11pt,a4paper]{article}

%====================================================
% PACKAGES
%====================================================
\usepackage[utf8]{inputenc}
\usepackage[T1]{fontenc}
\usepackage[french]{babel}
\usepackage[a4paper,margin=2cm]{geometry}
\usepackage{amsmath,amssymb,amsthm,mathtools}
\usepackage{enumitem}
\usepackage{fancyhdr}
\usepackage{lastpage}
\usepackage{xcolor}
\usepackage{array}
\usepackage{booktabs}

%====================================================
% MISE EN PAGE
%====================================================
\setlength{\headheight}{14pt}
\pagestyle{fancy}
\fancyhf{}
\fancyhead[L]{Classe préparatoire}
\fancyhead[C]{Corrigé -- DS Algèbre linéaire}
\fancyhead[R]{Page \thepage/\pageref{LastPage}}
\fancyfoot[C]{}

\setlength{\parindent}{0pt}
\setlength{\parskip}{0.45em}

\definecolor{bleu}{RGB}{20,60,120}
\definecolor{gris}{RGB}{85,85,85}

\newcommand{\pts}[1]{\hfill \textbf{[#1 pts]}}
\newcommand{\R}{\mathbb{R}}
\newcommand{\Ker}{\operatorname{Ker}}
\newcommand{\Img}{\operatorname{Im}}
\newcommand{\rg}{\operatorname{rg}}
\newcommand{\Mat}{\operatorname{Mat}}
\newcommand{\Vect}{\operatorname{Vect}}
\newcommand{\Id}{\operatorname{Id}}

\newcounter{probleme}
\newenvironment{probleme}[2]{
    \refstepcounter{probleme}
    \vspace{0.6cm}
    {\Large\color{bleu}\textbf{Problème \theprobleme{} -- #1}}\hfill \textbf{#2 pts}
    \par\medskip
    \hrule
    \medskip
}{\vspace{0.5cm}}

\newcommand{\partie}[1]{
    \vspace{0.35cm}
    {\large\color{gris}\textbf{#1}}
    \par\smallskip
}

\newcommand{\solution}{
    \textbf{\color{bleu}Solution. }
}

\setlist[enumerate,1]{label=\textbf{\arabic*.}, leftmargin=1.1cm}
\setlist[enumerate,2]{label=\textbf{\alph*)}, leftmargin=1cm}
\setlist[itemize]{leftmargin=1cm}

%====================================================
% DOCUMENT
%====================================================
\begin{document}

\begin{center}
    {\Huge \textbf{Corrigé détaillé}}\\[0.2cm]
    {\Large \textbf{Devoir surveillé -- Algèbre linéaire}}\\[0.2cm]
    {\large Systèmes linéaires, matrices, applications linéaires, rang et déterminants}\\[0.4cm]
\end{center}

\begin{center}
\begin{tabular}{|>{\bfseries}c|c|}
\hline
Niveau & Bac + 1 -- classe préparatoire scientifique \\
\hline
Total & 85 points \\
\hline
\end{tabular}
\end{center}

\vspace{0.3cm}

\textbf{Remarque générale.}  
Ce corrigé est rédigé dans un esprit concours : les calculs sont détaillés, les implications importantes sont justifiées, et les conclusions sont formulées explicitement.

%====================================================
% PROBLEME 1
%====================================================
\begin{probleme}{Un système linéaire à paramètre}{28}

On considère, pour un réel \(m\), le système linéaire \((S_m)\) d'inconnues \(x,y,z\) :
\[
(S_m)\qquad
\begin{cases}
x+y+z=1,\\
x+my+z=2,\\
x+y+mz=3.
\end{cases}
\]

On note
\[
A_m=
\begin{pmatrix}
1&1&1\\
1&m&1\\
1&1&m
\end{pmatrix},
\qquad
X=
\begin{pmatrix}
x\\y\\z
\end{pmatrix},
\qquad
B=
\begin{pmatrix}
1\\2\\3
\end{pmatrix}.
\]

\partie{Partie A -- Étude classique du système}

\begin{enumerate}
    \item Écrire le système \((S_m)\) sous la forme matricielle
    \[
    A_mX=B.
    \]
    \pts{1}

    \solution

    On a :
    \[
    A_mX=
    \begin{pmatrix}
    1&1&1\\
    1&m&1\\
    1&1&m
    \end{pmatrix}
    \begin{pmatrix}
    x\\y\\z
    \end{pmatrix}
    =
    \begin{pmatrix}
    x+y+z\\
    x+my+z\\
    x+y+mz
    \end{pmatrix}.
    \]
    Ainsi le système \((S_m)\) s'écrit bien :
    \[
    A_mX=
    \begin{pmatrix}
    1\\2\\3
    \end{pmatrix}
    =B.
    \]

    \item Calculer le déterminant de \(A_m\). \pts{3}

    \solution

    On calcule :
    \[
    \det(A_m)=
    \begin{vmatrix}
    1&1&1\\
    1&m&1\\
    1&1&m
    \end{vmatrix}.
    \]
    On effectue les opérations élémentaires sur les lignes :
    \[
    L_2\leftarrow L_2-L_1,
    \qquad
    L_3\leftarrow L_3-L_1.
    \]
    Ces opérations ne changent pas le déterminant. On obtient :
    \[
    \det(A_m)=
    \begin{vmatrix}
    1&1&1\\
    0&m-1&0\\
    0&0&m-1
    \end{vmatrix}.
    \]
    Cette matrice est triangulaire supérieure, donc son déterminant est le produit de ses coefficients diagonaux :
    \[
    \det(A_m)=1\cdot (m-1)\cdot (m-1).
    \]
    Finalement :
    \[
    \boxed{\det(A_m)=(m-1)^2.}
    \]

    \item Pour quelles valeurs de \(m\) la matrice \(A_m\) est-elle inversible ? \pts{2}

    \solution

    Une matrice carrée est inversible si et seulement si son déterminant est non nul. Or :
    \[
    \det(A_m)=(m-1)^2.
    \]
    Donc :
    \[
    \det(A_m)\neq 0
    \Longleftrightarrow
    (m-1)^2\neq 0
    \Longleftrightarrow
    m\neq 1.
    \]
    Ainsi :
    \[
    \boxed{A_m \text{ est inversible si et seulement si } m\neq 1.}
    \]

    \item On suppose dans cette question que \(m\neq 1\).
    Résoudre le système \((S_m)\). \pts{5}

    \solution

    On part du système :
    \[
    \begin{cases}
    x+y+z=1,\\
    x+my+z=2,\\
    x+y+mz=3.
    \end{cases}
    \]
    On soustrait la première équation à la deuxième :
    \[
    (x+my+z)-(x+y+z)=2-1.
    \]
    Donc :
    \[
    (m-1)y=1.
    \]
    Comme \(m\neq 1\), on peut diviser par \(m-1\), d'où :
    \[
    y=\frac{1}{m-1}.
    \]

    De même, on soustrait la première équation à la troisième :
    \[
    (x+y+mz)-(x+y+z)=3-1.
    \]
    Donc :
    \[
    (m-1)z=2.
    \]
    Comme \(m\neq 1\), on obtient :
    \[
    z=\frac{2}{m-1}.
    \]

    Il reste à utiliser la première équation :
    \[
    x+y+z=1.
    \]
    Ainsi :
    \[
    x=1-y-z
    =
    1-\frac{1}{m-1}-\frac{2}{m-1}
    =
    1-\frac{3}{m-1}.
    \]
    Donc :
    \[
    x=\frac{m-1}{m-1}-\frac{3}{m-1}
    =
    \frac{m-4}{m-1}.
    \]

    Finalement, pour \(m\neq 1\), l'unique solution est :
    \[
    \boxed{
    (x,y,z)
    =
    \left(
    \frac{m-4}{m-1},
    \frac{1}{m-1},
    \frac{2}{m-1}
    \right).}
    \]

    \item Étudier séparément le cas \(m=1\).
    Le système est-il compatible ? Si oui, donner l'ensemble de ses solutions. \pts{3}

    \solution

    Lorsque \(m=1\), le système devient :
    \[
    \begin{cases}
    x+y+z=1,\\
    x+y+z=2,\\
    x+y+z=3.
    \end{cases}
    \]
    Les trois équations imposent simultanément :
    \[
    x+y+z=1,\qquad x+y+z=2,\qquad x+y+z=3.
    \]
    C'est impossible, car un même réel ne peut pas être égal à \(1\), \(2\) et \(3\) en même temps.

    Donc le système est incompatible.

    Ainsi :
    \[
    \boxed{\text{Pour }m=1,\ (S_m)\text{ n'admet aucune solution.}}
    \]

    \item En déduire, selon les valeurs de \(m\), si le système \((S_m)\) admet :
    \begin{itemize}
        \item une unique solution ;
        \item aucune solution ;
        \item une infinité de solutions.
    \end{itemize}
    \pts{2}

    \solution

    D'après les questions précédentes :
    \begin{itemize}
        \item si \(m\neq 1\), la matrice \(A_m\) est inversible, donc le système admet une unique solution ;
        \item si \(m=1\), le système est incompatible, donc il n'admet aucune solution.
    \end{itemize}

    Il n'y a donc jamais une infinité de solutions pour le second membre \(B=(1,2,3)^T\).

    En conclusion :
    \[
    \boxed{
    \begin{cases}
    m\neq 1 : \text{ une unique solution},\\
    m=1 : \text{ aucune solution}.
    \end{cases}}
    \]
\end{enumerate}

\partie{Partie B -- Rang et système homogène associé}

On considère le système homogène associé :
\[
(H_m)\qquad A_mX=0.
\]

\begin{enumerate}[resume]
    \item Déterminer le rang de \(A_m\) selon les valeurs de \(m\). \pts{3}

    \solution

    On sait que :
    \[
    \det(A_m)=(m-1)^2.
    \]

    Si \(m\neq 1\), alors \(\det(A_m)\neq 0\). La matrice \(A_m\) est donc inversible, et son rang est maximal :
    \[
    \rg(A_m)=3.
    \]

    Si \(m=1\), alors :
    \[
    A_1=
    \begin{pmatrix}
    1&1&1\\
    1&1&1\\
    1&1&1
    \end{pmatrix}.
    \]
    Toutes les lignes sont égales et non nulles. La matrice possède donc une seule direction de lignes non nulle. Ainsi :
    \[
    \rg(A_1)=1.
    \]

    Finalement :
    \[
    \boxed{
    \rg(A_m)=
    \begin{cases}
    3 & \text{si } m\neq 1,\\
    1 & \text{si } m=1.
    \end{cases}}
    \]

    \item Résoudre le système homogène \((H_m)\) selon les valeurs de \(m\). \pts{3}

    \solution

    Le système homogène s'écrit :
    \[
    A_mX=0.
    \]

    Si \(m\neq 1\), la matrice \(A_m\) est inversible. Donc :
    \[
    A_mX=0
    \Longrightarrow
    X=A_m^{-1}0=0.
    \]
    Ainsi :
    \[
    \boxed{\Ker(A_m)=\{(0,0,0)\}\quad\text{si }m\neq 1.}
    \]

    Si \(m=1\), le système devient :
    \[
    \begin{cases}
    x+y+z=0,\\
    x+y+z=0,\\
    x+y+z=0.
    \end{cases}
    \]
    Il est donc équivalent à :
    \[
    x+y+z=0.
    \]
    On peut poser \(y=s\) et \(z=t\), avec \(s,t\in\R\). Alors :
    \[
    x=-s-t.
    \]
    Donc :
    \[
    (x,y,z)=(-s-t,s,t)
    =
    s(-1,1,0)+t(-1,0,1).
    \]
    Ainsi :
    \[
    \boxed{
    \Ker(A_1)=
    \Vect\big((-1,1,0),(-1,0,1)\big).}
    \]

    \item Expliquer pourquoi le résultat obtenu est cohérent avec le théorème du rang appliqué à l'application linéaire
    \[
    u_m:\R^3\longrightarrow\R^3,
    \qquad
    X\longmapsto A_mX.
    \]
    \pts{3}

    \solution

    L'application \(u_m\) est une application linéaire de \(\R^3\) dans \(\R^3\). Son noyau est exactement l'ensemble des solutions du système homogène :
    \[
    \Ker(u_m)=\{X\in\R^3\mid A_mX=0\}.
    \]

    Le théorème du rang donne :
    \[
    \dim(\R^3)=\dim(\Ker(u_m))+\rg(u_m).
    \]
    Comme la matrice de \(u_m\) est \(A_m\), on a :
    \[
    \rg(u_m)=\rg(A_m).
    \]

    Si \(m\neq 1\), on a \(\rg(A_m)=3\). Donc :
    \[
    3=\dim(\Ker(u_m))+3,
    \]
    d'où :
    \[
    \dim(\Ker(u_m))=0.
    \]
    Cela correspond bien au noyau réduit à \(\{0\}\).

    Si \(m=1\), on a \(\rg(A_1)=1\). Donc :
    \[
    3=\dim(\Ker(u_1))+1,
    \]
    d'où :
    \[
    \dim(\Ker(u_1))=2.
    \]
    Cela correspond bien au noyau trouvé :
    \[
    \Ker(A_1)=
    \Vect\big((-1,1,0),(-1,0,1)\big),
    \]
    qui est un plan vectoriel de dimension \(2\).

    Les résultats sont donc cohérents avec le théorème du rang.
\end{enumerate}

\partie{Partie C -- Généralisation}

On remplace maintenant le second membre \(B\) par un second membre quelconque
\[
Y=
\begin{pmatrix}
a\\b\\c
\end{pmatrix}
\in\R^3.
\]
On considère donc le système
\[
A_mX=Y.
\]

\begin{enumerate}[resume]
    \item Lorsque \(m\neq 1\), justifier que le système admet une unique solution pour tout \(Y\in\R^3\). \pts{2}

    \solution

    Si \(m\neq 1\), on a :
    \[
    \det(A_m)=(m-1)^2\neq 0.
    \]
    Donc \(A_m\) est inversible.

    Ainsi, pour tout second membre \(Y\in\R^3\), le système
    \[
    A_mX=Y
    \]
    admet une unique solution donnée par :
    \[
    X=A_m^{-1}Y.
    \]

    Donc :
    \[
    \boxed{\text{Si }m\neq 1,\text{ le système }A_mX=Y\text{ admet une unique solution pour tout }Y\in\R^3.}
    \]

    \item On suppose maintenant \(m=1\). Donner une condition nécessaire et suffisante portant sur \(a,b,c\) pour que le système
    \[
    A_1X=Y
    \]
    soit compatible. \pts{3}

    \solution

    Lorsque \(m=1\), on a :
    \[
    A_1=
    \begin{pmatrix}
    1&1&1\\
    1&1&1\\
    1&1&1
    \end{pmatrix}.
    \]
    Le système
    \[
    A_1X=Y
    \]
    s'écrit :
    \[
    \begin{cases}
    x+y+z=a,\\
    x+y+z=b,\\
    x+y+z=c.
    \end{cases}
    \]
    Pour que ce système soit compatible, il faut et il suffit que les trois seconds membres soient égaux :
    \[
    a=b=c.
    \]

    En effet, si \(a=b=c\), le système se réduit à une seule équation :
    \[
    x+y+z=a,
    \]
    qui admet des solutions.

    Donc :
    \[
    \boxed{A_1X=Y\text{ est compatible si et seulement si }a=b=c.}
    \]

    \item Dans le cas \(m=1\), interpréter géométriquement l'ensemble des seconds membres \(Y\) pour lesquels le système est compatible. \pts{1}

    \solution

    Les seconds membres compatibles sont les vecteurs :
    \[
    Y=
    \begin{pmatrix}
    a\\a\\a
    \end{pmatrix}
    =
    a
    \begin{pmatrix}
    1\\1\\1
    \end{pmatrix},
    \qquad a\in\R.
    \]
    Ils forment donc la droite vectorielle :
    \[
    \boxed{\Vect\left(
    \begin{pmatrix}
    1\\1\\1
    \end{pmatrix}
    \right).}
    \]
    Géométriquement, il s'agit d'une droite vectorielle de \(\R^3\).
\end{enumerate}

\end{probleme}

%====================================================
% PROBLEME 2
%====================================================
\begin{probleme}{La machine à mélanger les jetons}{27}

Une petite machine reçoit trois nombres réels \(x,y,z\), représentant les quantités de trois types de jetons : rouges, bleus et verts. Elle renvoie trois nouvelles quantités \(a,b,c\) selon la règle :
\[
\begin{cases}
a=x+y,\\
b=y+z,\\
c=x+z.
\end{cases}
\]

On définit ainsi une application
\[
T:\R^3\longrightarrow \R^3,
\qquad
T(x,y,z)=(x+y,\ y+z,\ x+z).
\]

\partie{Partie A -- Étude de la machine complète}

\begin{enumerate}
    \item Montrer que \(T\) est une application linéaire. \pts{2}

    \solution

    Soient \(X=(x,y,z)\) et \(X'=(x',y',z')\) deux vecteurs de \(\R^3\), et soient \(\lambda,\mu\in\R\).

    On a :
    \[
    \lambda X+\mu X'
    =
    (\lambda x+\mu x',\lambda y+\mu y',\lambda z+\mu z').
    \]
    Alors :
    \[
    T(\lambda X+\mu X')
    =
    \big(
    \lambda x+\mu x'+\lambda y+\mu y',
    \lambda y+\mu y'+\lambda z+\mu z',
    \lambda x+\mu x'+\lambda z+\mu z'
    \big).
    \]
    En regroupant les termes :
    \[
    T(\lambda X+\mu X')
    =
    \lambda(x+y,y+z,x+z)
    +
    \mu(x'+y',y'+z',x'+z').
    \]
    Donc :
    \[
    T(\lambda X+\mu X')
    =
    \lambda T(X)+\mu T(X').
    \]
    Ainsi :
    \[
    \boxed{T\text{ est linéaire}.}
    \]

    \item Donner la matrice \(M\) de \(T\) dans la base canonique de \(\R^3\). \pts{2}

    \solution

    On calcule les images des vecteurs de la base canonique :
    \[
    e_1=(1,0,0),\quad e_2=(0,1,0),\quad e_3=(0,0,1).
    \]
    On obtient :
    \[
    T(e_1)=T(1,0,0)=(1,0,1),
    \]
    \[
    T(e_2)=T(0,1,0)=(1,1,0),
    \]
    \[
    T(e_3)=T(0,0,1)=(0,1,1).
    \]
    La matrice de \(T\) dans la base canonique a pour colonnes ces trois vecteurs :
    \[
    \boxed{
    M=
    \begin{pmatrix}
    1&1&0\\
    0&1&1\\
    1&0&1
    \end{pmatrix}.}
    \]

    \item Calculer le déterminant de \(M\). \pts{3}

    \solution

    On calcule :
    \[
    \det(M)=
    \begin{vmatrix}
    1&1&0\\
    0&1&1\\
    1&0&1
    \end{vmatrix}.
    \]
    En développant selon la première ligne :
    \[
    \det(M)
    =
    1
    \begin{vmatrix}
    1&1\\
    0&1
    \end{vmatrix}
    -
    1
    \begin{vmatrix}
    0&1\\
    1&1
    \end{vmatrix}
    +
    0.
    \]
    Donc :
    \[
    \det(M)
    =
    1(1\cdot 1-1\cdot 0)
    -
    1(0\cdot 1-1\cdot 1).
    \]
    Ainsi :
    \[
    \det(M)=1-(-1)=2.
    \]
    Donc :
    \[
    \boxed{\det(M)=2.}
    \]

    \item En déduire que \(T\) est bijective. \pts{2}

    \solution

    La matrice de \(T\) dans la base canonique est \(M\). Or :
    \[
    \det(M)=2\neq 0.
    \]
    Donc \(M\) est inversible.

    Une application linéaire de \(\R^3\) dans \(\R^3\) dont la matrice est inversible est un automorphisme. Par conséquent :
    \[
    \boxed{T\text{ est bijective}.}
    \]

    \item La machine affiche en sortie
    \[
    (a,b,c)=(7,8,9).
    \]
    Déterminer les valeurs initiales \(x,y,z\). \pts{3}

    \solution

    On doit résoudre :
    \[
    T(x,y,z)=(7,8,9).
    \]
    Cela donne le système :
    \[
    \begin{cases}
    x+y=7,\\
    y+z=8,\\
    x+z=9.
    \end{cases}
    \]
    Additionnons la première et la troisième équation :
    \[
    (x+y)+(x+z)=7+9.
    \]
    Donc :
    \[
    2x+y+z=16.
    \]
    Or, d'après la deuxième équation :
    \[
    y+z=8.
    \]
    Ainsi :
    \[
    2x+8=16,
    \]
    donc :
    \[
    x=4.
    \]
    Puis :
    \[
    x+y=7
    \Longrightarrow
    4+y=7
    \Longrightarrow
    y=3.
    \]
    Enfin :
    \[
    y+z=8
    \Longrightarrow
    3+z=8
    \Longrightarrow
    z=5.
    \]
    Donc :
    \[
    \boxed{(x,y,z)=(4,3,5).}
    \]

    \item Résoudre plus généralement le système
    \[
    T(x,y,z)=(a,b,c),
    \]
    où \(a,b,c\) sont trois réels quelconques. \pts{4}

    \solution

    On résout :
    \[
    \begin{cases}
    x+y=a,\\
    y+z=b,\\
    x+z=c.
    \end{cases}
    \]
    On additionne la première et la troisième équation :
    \[
    (x+y)+(x+z)=a+c.
    \]
    Donc :
    \[
    2x+y+z=a+c.
    \]
    Or :
    \[
    y+z=b.
    \]
    Ainsi :
    \[
    2x+b=a+c,
    \]
    d'où :
    \[
    x=\frac{a-b+c}{2}.
    \]

    De même, en additionnant les deux premières équations :
    \[
    x+2y+z=a+b.
    \]
    Comme \(x+z=c\), on obtient :
    \[
    c+2y=a+b,
    \]
    d'où :
    \[
    y=\frac{a+b-c}{2}.
    \]

    Enfin, en additionnant les deux dernières équations :
    \[
    x+y+2z=b+c.
    \]
    Comme \(x+y=a\), on obtient :
    \[
    a+2z=b+c,
    \]
    d'où :
    \[
    z=\frac{-a+b+c}{2}.
    \]

    Finalement :
    \[
    \boxed{
    (x,y,z)=
    \left(
    \frac{a-b+c}{2},
    \frac{a+b-c}{2},
    \frac{-a+b+c}{2}
    \right).}
    \]

    \item En déduire la matrice \(M^{-1}\). \pts{3}

    \solution

    D'après la question précédente, si :
    \[
    \begin{pmatrix}
    a\\b\\c
    \end{pmatrix}
    =
    M
    \begin{pmatrix}
    x\\y\\z
    \end{pmatrix},
    \]
    alors :
    \[
    \begin{pmatrix}
    x\\y\\z
    \end{pmatrix}
    =
    \begin{pmatrix}
    \dfrac{a-b+c}{2}\\[0.2cm]
    \dfrac{a+b-c}{2}\\[0.2cm]
    \dfrac{-a+b+c}{2}
    \end{pmatrix}.
    \]
    Donc :
    \[
    \begin{pmatrix}
    x\\y\\z
    \end{pmatrix}
    =
    \frac12
    \begin{pmatrix}
    1&-1&1\\
    1&1&-1\\
    -1&1&1
    \end{pmatrix}
    \begin{pmatrix}
    a\\b\\c
    \end{pmatrix}.
    \]
    Par conséquent :
    \[
    \boxed{
    M^{-1}
    =
    \frac12
    \begin{pmatrix}
    1&-1&1\\
    1&1&-1\\
    -1&1&1
    \end{pmatrix}.}
    \]
\end{enumerate}

\partie{Partie B -- Machine avec perte d'information}

On modifie maintenant la machine. Elle renvoie seulement les deux nombres
\[
u=x+y,
\qquad
v=y+z.
\]
On définit donc
\[
F:\R^3\longrightarrow \R^2,
\qquad
F(x,y,z)=(x+y,\ y+z).
\]

\begin{enumerate}[resume]
    \item Donner la matrice \(N\) de \(F\) dans les bases canoniques. \pts{1}

    \solution

    On a :
    \[
    F(x,y,z)=(x+y,y+z).
    \]
    Donc :
    \[
    F(e_1)=F(1,0,0)=(1,0),
    \]
    \[
    F(e_2)=F(0,1,0)=(1,1),
    \]
    \[
    F(e_3)=F(0,0,1)=(0,1).
    \]
    La matrice de \(F\) dans les bases canoniques est donc :
    \[
    \boxed{
    N=
    \begin{pmatrix}
    1&1&0\\
    0&1&1
    \end{pmatrix}.}
    \]

    \item Déterminer \(\Ker(F)\). \pts{3}

    \solution

    On cherche les vecteurs \((x,y,z)\in\R^3\) tels que :
    \[
    F(x,y,z)=(0,0).
    \]
    Cela équivaut à :
    \[
    \begin{cases}
    x+y=0,\\
    y+z=0.
    \end{cases}
    \]
    De la première équation :
    \[
    x=-y.
    \]
    De la deuxième :
    \[
    z=-y.
    \]
    En posant \(y=t\), avec \(t\in\R\), on obtient :
    \[
    (x,y,z)=(-t,t,-t)=t(-1,1,-1).
    \]
    Donc :
    \[
    \boxed{\Ker(F)=\Vect((-1,1,-1)).}
    \]

    \item En déduire le rang de \(F\) à l'aide du théorème du rang. \pts{2}

    \solution

    L'application \(F\) est une application linéaire de \(\R^3\) dans \(\R^2\). Le théorème du rang donne :
    \[
    \dim(\R^3)=\dim(\Ker(F))+\rg(F).
    \]
    Or :
    \[
    \Ker(F)=\Vect((-1,1,-1)),
    \]
    donc :
    \[
    \dim(\Ker(F))=1.
    \]
    Ainsi :
    \[
    3=1+\rg(F),
    \]
    d'où :
    \[
    \boxed{\rg(F)=2.}
    \]

    \item L'application \(F\) est-elle injective ? Est-elle surjective ? Justifier. \pts{3}

    \solution

    Une application linéaire est injective si et seulement si son noyau est réduit au vecteur nul.

    Ici :
    \[
    \Ker(F)=\Vect((-1,1,-1)).
    \]
    Ce noyau n'est pas réduit à \(\{0\}\). Donc :
    \[
    \boxed{F\text{ n'est pas injective}.}
    \]

    Ensuite, \(F\) est une application de \(\R^3\) dans \(\R^2\). On a montré que :
    \[
    \rg(F)=2.
    \]
    Or :
    \[
    \dim(\R^2)=2.
    \]
    Donc :
    \[
    \rg(F)=\dim(\R^2).
    \]
    L'image de \(F\) est donc tout \(\R^2\). Ainsi :
    \[
    \boxed{F\text{ est surjective}.}
    \]
\end{enumerate}

\partie{Partie C -- Une généralisation à paramètre}

Pour un réel \(\lambda\), on définit une nouvelle machine
\[
T_\lambda:\R^3\longrightarrow\R^3
\]
par
\[
T_\lambda(x,y,z)=(x+y,\ y+z,\ x+\lambda z).
\]

\begin{enumerate}[resume]
    \item Donner la matrice \(M_\lambda\) de \(T_\lambda\) dans la base canonique. \pts{1}

    \solution

    On lit directement :
    \[
    T_\lambda(x,y,z)
    =
    (x+y,\ y+z,\ x+\lambda z).
    \]
    La matrice dans la base canonique est :
    \[
    \boxed{
    M_\lambda=
    \begin{pmatrix}
    1&1&0\\
    0&1&1\\
    1&0&\lambda
    \end{pmatrix}.}
    \]

    \item Calculer \(\det(M_\lambda)\). \pts{2}

    \solution

    On calcule :
    \[
    \det(M_\lambda)=
    \begin{vmatrix}
    1&1&0\\
    0&1&1\\
    1&0&\lambda
    \end{vmatrix}.
    \]
    En développant selon la première ligne :
    \[
    \det(M_\lambda)
    =
    1
    \begin{vmatrix}
    1&1\\
    0&\lambda
    \end{vmatrix}
    -
    1
    \begin{vmatrix}
    0&1\\
    1&\lambda
    \end{vmatrix}.
    \]
    Donc :
    \[
    \det(M_\lambda)
    =
    1(1\cdot\lambda-1\cdot 0)
    -
    1(0\cdot\lambda-1\cdot 1).
    \]
    Ainsi :
    \[
    \det(M_\lambda)=\lambda-(-1)=\lambda+1.
    \]
    Donc :
    \[
    \boxed{\det(M_\lambda)=\lambda+1.}
    \]

    \item Pour quelles valeurs de \(\lambda\) la machine \(T_\lambda\) permet-elle toujours de retrouver les valeurs initiales \(x,y,z\) à partir de la sortie ? \pts{2}

    \solution

    Retrouver toujours les valeurs initiales à partir de la sortie signifie que l'application \(T_\lambda\) est bijective.

    Comme \(T_\lambda\) est une application linéaire de \(\R^3\) dans \(\R^3\), elle est bijective si et seulement si sa matrice est inversible, c'est-à-dire si et seulement si :
    \[
    \det(M_\lambda)\neq 0.
    \]
    Or :
    \[
    \det(M_\lambda)=\lambda+1.
    \]
    Donc :
    \[
    \lambda+1\neq 0
    \Longleftrightarrow
    \lambda\neq -1.
    \]

    Ainsi :
    \[
    \boxed{T_\lambda\text{ permet toujours de retrouver les valeurs initiales si et seulement si }\lambda\neq -1.}
    \]
\end{enumerate}

\end{probleme}

%====================================================
% PROBLEME 3
%====================================================
\begin{probleme}{Une application linéaire presque symétrique}{30}

On considère l'application
\[
u:\R^3\longrightarrow \R^3
\]
définie par
\[
u(x,y,z)=(x+y,\ y+z,\ x+z).
\]

On note \(B\) la base canonique de \(\R^3\).

\partie{Partie A -- Une application bijective}

\begin{enumerate}
    \item Justifier que \(u\) est linéaire. \pts{1}

    \solution

    L'application \(u\) est définie par :
    \[
    u(x,y,z)=(x+y,y+z,x+z).
    \]
    Chacune des coordonnées de \(u(x,y,z)\) est une combinaison linéaire des coordonnées \(x,y,z\). Donc \(u\) est linéaire.

    Plus explicitement, pour tous \(X,X'\in\R^3\) et tous \(\lambda,\mu\in\R\), on vérifie :
    \[
    u(\lambda X+\mu X')=\lambda u(X)+\mu u(X').
    \]
    Ainsi :
    \[
    \boxed{u\text{ est linéaire}.}
    \]

    \item Donner la matrice \(A=\Mat_B(u)\). \pts{2}

    \solution

    On calcule les images des vecteurs de la base canonique :
    \[
    u(1,0,0)=(1,0,1),
    \]
    \[
    u(0,1,0)=(1,1,0),
    \]
    \[
    u(0,0,1)=(0,1,1).
    \]
    Donc :
    \[
    \boxed{
    A=
    \begin{pmatrix}
    1&1&0\\
    0&1&1\\
    1&0&1
    \end{pmatrix}.}
    \]

    \item Calculer \(\det(A)\). \pts{3}

    \solution

    On reconnaît la même matrice que dans le problème précédent :
    \[
    A=
    \begin{pmatrix}
    1&1&0\\
    0&1&1\\
    1&0&1
    \end{pmatrix}.
    \]
    On calcule :
    \[
    \det(A)=
    \begin{vmatrix}
    1&1&0\\
    0&1&1\\
    1&0&1
    \end{vmatrix}.
    \]
    En développant selon la première ligne :
    \[
    \det(A)
    =
    1
    \begin{vmatrix}
    1&1\\
    0&1
    \end{vmatrix}
    -
    1
    \begin{vmatrix}
    0&1\\
    1&1
    \end{vmatrix}.
    \]
    Donc :
    \[
    \det(A)=1-(-1)=2.
    \]
    Ainsi :
    \[
    \boxed{\det(A)=2.}
    \]

    \item En déduire que \(u\) est un automorphisme de \(\R^3\). \pts{2}

    \solution

    Comme :
    \[
    \det(A)=2\neq 0,
    \]
    la matrice \(A\) est inversible.

    L'application \(u\) est une application linéaire de \(\R^3\) dans \(\R^3\), et sa matrice dans la base canonique est inversible. Donc \(u\) est bijective.

    Ainsi :
    \[
    \boxed{u\text{ est un automorphisme de }\R^3.}
    \]

    \item Calculer explicitement \(\Ker(u)\). Vérifier que le résultat est cohérent avec la question précédente. \pts{3}

    \solution

    On cherche les vecteurs \((x,y,z)\in\R^3\) tels que :
    \[
    u(x,y,z)=(0,0,0).
    \]
    Cela donne :
    \[
    \begin{cases}
    x+y=0,\\
    y+z=0,\\
    x+z=0.
    \end{cases}
    \]
    De la première équation :
    \[
    x=-y.
    \]
    De la deuxième :
    \[
    z=-y.
    \]
    En remplaçant dans la troisième :
    \[
    x+z=-y-y=-2y=0.
    \]
    Donc :
    \[
    y=0.
    \]
    Par conséquent :
    \[
    x=0,\qquad z=0.
    \]
    Ainsi :
    \[
    \boxed{\Ker(u)=\{(0,0,0)\}.}
    \]

    Ce résultat est cohérent avec le fait que \(u\) est un automorphisme : une application linéaire est injective si et seulement si son noyau est réduit au vecteur nul.

    \item Déterminer \(\Img(u)\). \pts{2}

    \solution

    Puisque \(u\) est un automorphisme de \(\R^3\), elle est en particulier surjective. Donc :
    \[
    \Img(u)=\R^3.
    \]
    Ainsi :
    \[
    \boxed{\Img(u)=\R^3.}
    \]
\end{enumerate}

\partie{Partie B -- Une application proche mais non bijective}

On considère maintenant l'application
\[
v:\R^3\longrightarrow \R^3
\]
définie par
\[
v(x,y,z)=(x+y,\ y+z,\ x+2y+z).
\]

\begin{enumerate}[resume]
    \item Montrer que \(v\) est linéaire et donner sa matrice \(C\) dans la base canonique. \pts{2}

    \solution

    Les coordonnées de \(v(x,y,z)\) sont :
    \[
    x+y,\qquad y+z,\qquad x+2y+z.
    \]
    Ce sont des combinaisons linéaires de \(x,y,z\). Donc \(v\) est linéaire.

    Sa matrice dans la base canonique est obtenue par lecture des coefficients :
    \[
    \boxed{
    C=
    \begin{pmatrix}
    1&1&0\\
    0&1&1\\
    1&2&1
    \end{pmatrix}.}
    \]

    \item Calculer \(\det(C)\). \pts{2}

    \solution

    On remarque que la troisième ligne est la somme des deux premières :
    \[
    (1,2,1)=(1,1,0)+(0,1,1).
    \]
    Les lignes de \(C\) sont donc liées. Par conséquent :
    \[
    \boxed{\det(C)=0.}
    \]

    On pouvait également le vérifier par calcul direct.

    \item Déterminer \(\Ker(v)\). \pts{3}

    \solution

    On cherche \((x,y,z)\in\R^3\) tel que :
    \[
    v(x,y,z)=(0,0,0).
    \]
    Cela donne :
    \[
    \begin{cases}
    x+y=0,\\
    y+z=0,\\
    x+2y+z=0.
    \end{cases}
    \]
    La troisième équation est la somme des deux premières :
    \[
    (x+y)+(y+z)=x+2y+z.
    \]
    Elle est donc automatiquement vérifiée si les deux premières le sont.

    On résout alors :
    \[
    x+y=0,
    \qquad
    y+z=0.
    \]
    D'où :
    \[
    x=-y,\qquad z=-y.
    \]
    En posant \(y=t\), on obtient :
    \[
    (x,y,z)=(-t,t,-t)=t(-1,1,-1).
    \]
    Donc :
    \[
    \boxed{\Ker(v)=\Vect((-1,1,-1)).}
    \]

    \item En déduire le rang de \(v\) à l'aide du théorème du rang. \pts{2}

    \solution

    L'application \(v\) est linéaire de \(\R^3\) dans \(\R^3\). Le théorème du rang donne :
    \[
    \dim(\R^3)=\dim(\Ker(v))+\rg(v).
    \]
    Or :
    \[
    \Ker(v)=\Vect((-1,1,-1)),
    \]
    donc :
    \[
    \dim(\Ker(v))=1.
    \]
    Ainsi :
    \[
    3=1+\rg(v),
    \]
    d'où :
    \[
    \boxed{\rg(v)=2.}
    \]

    \item Donner une base de \(\Img(v)\). \pts{3}

    \solution

    Les colonnes de la matrice \(C\) engendrent l'image de \(v\). On note :
    \[
    C_1=
    \begin{pmatrix}
    1\\0\\1
    \end{pmatrix},
    \qquad
    C_2=
    \begin{pmatrix}
    1\\1\\2
    \end{pmatrix},
    \qquad
    C_3=
    \begin{pmatrix}
    0\\1\\1
    \end{pmatrix}.
    \]
    On remarque que :
    \[
    C_2=C_1+C_3.
    \]
    Ainsi :
    \[
    \Img(v)=\Vect(C_1,C_2,C_3)=\Vect(C_1,C_3).
    \]
    Les vecteurs \(C_1\) et \(C_3\) ne sont pas colinéaires. Ils forment donc une famille libre.

    Comme \(\rg(v)=2\), une famille libre de deux vecteurs de \(\Img(v)\) est une base de \(\Img(v)\).

    Ainsi :
    \[
    \boxed{
    \left(
    \begin{pmatrix}
    1\\0\\1
    \end{pmatrix},
    \begin{pmatrix}
    0\\1\\1
    \end{pmatrix}
    \right)
    \text{ est une base de }\Img(v).}
    \]

    \item Expliquer clairement pourquoi \(v\) n'est pas bijective. \pts{2}

    \solution

    Plusieurs arguments sont possibles.

    D'abord :
    \[
    \det(C)=0.
    \]
    Donc la matrice de \(v\) n'est pas inversible, et par conséquent \(v\) n'est pas bijective.

    On peut aussi dire que :
    \[
    \Ker(v)=\Vect((-1,1,-1))\neq\{0\}.
    \]
    Donc \(v\) n'est pas injective, et ne peut donc pas être bijective.

    Enfin :
    \[
    \rg(v)=2<3.
    \]
    Donc \(v\) n'est pas surjective sur \(\R^3\).

    Ainsi :
    \[
    \boxed{v\text{ n'est pas bijective}.}
    \]
\end{enumerate}

\partie{Partie C -- Une généralisation élégante}

Pour un réel \(\alpha\), on définit
\[
w_\alpha:\R^3\longrightarrow\R^3
\]
par
\[
w_\alpha(x,y,z)=(x+y,\ y+z,\ x+\alpha y+z).
\]

\begin{enumerate}[resume]
    \item Donner la matrice \(D_\alpha\) de \(w_\alpha\) dans la base canonique. \pts{1}

    \solution

    En lisant les coefficients de \(x,y,z\) dans chaque coordonnée :
    \[
    w_\alpha(x,y,z)=(x+y,\ y+z,\ x+\alpha y+z).
    \]
    La matrice dans la base canonique est :
    \[
    \boxed{
    D_\alpha=
    \begin{pmatrix}
    1&1&0\\
    0&1&1\\
    1&\alpha&1
    \end{pmatrix}.}
    \]

    \item Calculer \(\det(D_\alpha)\). \pts{3}

    \solution

    On calcule :
    \[
    \det(D_\alpha)=
    \begin{vmatrix}
    1&1&0\\
    0&1&1\\
    1&\alpha&1
    \end{vmatrix}.
    \]
    En développant selon la première ligne :
    \[
    \det(D_\alpha)
    =
    1
    \begin{vmatrix}
    1&1\\
    \alpha&1
    \end{vmatrix}
    -
    1
    \begin{vmatrix}
    0&1\\
    1&1
    \end{vmatrix}.
    \]
    Donc :
    \[
    \det(D_\alpha)
    =
    1(1-\alpha)
    -
    1(0-1).
    \]
    Ainsi :
    \[
    \det(D_\alpha)=1-\alpha+1=2-\alpha.
    \]
    Donc :
    \[
    \boxed{\det(D_\alpha)=2-\alpha.}
    \]

    \item Pour quelles valeurs de \(\alpha\) l'application \(w_\alpha\) est-elle un automorphisme de \(\R^3\) ? \pts{2}

    \solution

    L'application \(w_\alpha\) est un endomorphisme de \(\R^3\). Elle est un automorphisme si et seulement si sa matrice \(D_\alpha\) est inversible, c'est-à-dire si et seulement si :
    \[
    \det(D_\alpha)\neq 0.
    \]
    Or :
    \[
    \det(D_\alpha)=2-\alpha.
    \]
    Donc :
    \[
    2-\alpha\neq 0
    \Longleftrightarrow
    \alpha\neq 2.
    \]

    Ainsi :
    \[
    \boxed{w_\alpha\text{ est un automorphisme de }\R^3\text{ si et seulement si }\alpha\neq 2.}
    \]

    \item Dans le cas où \(w_\alpha\) n'est pas un automorphisme, déterminer \(\Ker(w_\alpha)\). \pts{3}

    \solution

    D'après la question précédente, \(w_\alpha\) n'est pas un automorphisme uniquement lorsque :
    \[
    \alpha=2.
    \]
    Dans ce cas :
    \[
    w_2(x,y,z)=(x+y,\ y+z,\ x+2y+z).
    \]
    C'est exactement l'application \(v\) étudiée dans la partie B.

    On cherche :
    \[
    w_2(x,y,z)=(0,0,0).
    \]
    Cela donne :
    \[
    \begin{cases}
    x+y=0,\\
    y+z=0,\\
    x+2y+z=0.
    \end{cases}
    \]
    La troisième équation est la somme des deux premières. On obtient donc :
    \[
    x=-y,\qquad z=-y.
    \]
    En posant \(y=t\), on trouve :
    \[
    (x,y,z)=(-t,t,-t)=t(-1,1,-1).
    \]
    Ainsi :
    \[
    \boxed{\Ker(w_2)=\Vect((-1,1,-1)).}
    \]

    \item Dans ce même cas, donner le rang de \(w_\alpha\), puis une base de \(\Img(w_\alpha)\). \pts{3}

    \solution

    On se place toujours dans le cas singulier :
    \[
    \alpha=2.
    \]
    On vient de montrer que :
    \[
    \Ker(w_2)=\Vect((-1,1,-1)).
    \]
    Donc :
    \[
    \dim(\Ker(w_2))=1.
    \]
    Par le théorème du rang :
    \[
    \dim(\R^3)=\dim(\Ker(w_2))+\rg(w_2).
    \]
    Ainsi :
    \[
    3=1+\rg(w_2),
    \]
    d'où :
    \[
    \boxed{\rg(w_2)=2.}
    \]

    Les colonnes de la matrice
    \[
    D_2=
    \begin{pmatrix}
    1&1&0\\
    0&1&1\\
    1&2&1
    \end{pmatrix}
    \]
    sont :
    \[
    C_1=
    \begin{pmatrix}
    1\\0\\1
    \end{pmatrix},
    \qquad
    C_2=
    \begin{pmatrix}
    1\\1\\2
    \end{pmatrix},
    \qquad
    C_3=
    \begin{pmatrix}
    0\\1\\1
    \end{pmatrix}.
    \]
    Or :
    \[
    C_2=C_1+C_3.
    \]
    Les vecteurs \(C_1\) et \(C_3\) sont non colinéaires. Comme l'image est de dimension \(2\), ils forment une base de l'image.

    Donc :
    \[
    \boxed{
    \left(
    \begin{pmatrix}
    1\\0\\1
    \end{pmatrix},
    \begin{pmatrix}
    0\\1\\1
    \end{pmatrix}
    \right)
    \text{ est une base de }\Img(w_2).}
    \]

    \item Comparer les applications \(u\), \(v\) et \(w_\alpha\) du point de vue de l'inversibilité, du noyau et du rang.
    Rédiger une conclusion synthétique. \pts{2}

    \solution

    L'application \(u\) a pour matrice :
    \[
    A=
    \begin{pmatrix}
    1&1&0\\
    0&1&1\\
    1&0&1
    \end{pmatrix},
    \]
    et :
    \[
    \det(A)=2\neq 0.
    \]
    Donc \(u\) est un automorphisme de \(\R^3\). On a :
    \[
    \Ker(u)=\{0\},
    \qquad
    \rg(u)=3,
    \qquad
    \Img(u)=\R^3.
    \]

    L'application \(v\) correspond au cas particulier \(\alpha=2\) de \(w_\alpha\). Sa matrice est non inversible, car :
    \[
    \det(C)=0.
    \]
    On a :
    \[
    \Ker(v)=\Vect((-1,1,-1)),
    \qquad
    \rg(v)=2.
    \]
    Donc \(v\) n'est pas bijective.

    Enfin, pour \(w_\alpha\), on a :
    \[
    \det(D_\alpha)=2-\alpha.
    \]
    Donc :
    \[
    \boxed{
    \begin{cases}
    \alpha\neq 2 : w_\alpha \text{ est un automorphisme},\ \Ker(w_\alpha)=\{0\},\ \rg(w_\alpha)=3,\\[0.15cm]
    \alpha=2 : w_\alpha \text{ n'est pas un automorphisme},\ \Ker(w_\alpha)=\Vect((-1,1,-1)),\ \rg(w_\alpha)=2.
    \end{cases}}
    \]

    La perte d'inversibilité apparaît exactement lorsque le déterminant s'annule. Dans ce cas, une direction non nulle est envoyée sur le vecteur nul, ce qui fait apparaître un noyau non trivial et diminue le rang.
\end{enumerate}

\end{probleme}

\vfill

\begin{center}
    \textbf{Fin du corrigé}
\end{center}

\end{document}